
\documentclass[11pt]{article}
\usepackage[margin=1in]{geometry}
\usepackage[T1]{fontenc}
\usepackage[utf8]{inputenc}
\usepackage{lmodern}
\usepackage{amsmath}
\usepackage{amssymb}
\usepackage{graphicx}
\usepackage{booktabs}
\usepackage{float}
\usepackage{hyperref}
\hypersetup{hidelinks}
\setlength{\parskip}{0.5em}
\setlength{\parindent}{0pt}
\begin{document}
\sloppy
\section{Step 3: fNIRS Trigger Harmonization, Session Quality Control, and Basic Preprocessing}
\label{sec:step3_fnirs}
% Analyst note added 2026-03-20: PID006 and PID037 are globally excluded from downstream participant-level analyses. PID006 was excluded after Step 3 documented session 2025-11-12_002 as an implausibly short 10.420224 second recording with no paired PsychoPy log. PID037 was excluded because a behavioral log exists but no fNIRS session folder was found. Step 3 raw QC outputs remain for audit, but future downstream steps should honor materials/analysis_participant_exclusions.json.

\subsection{Objective}
The objective of Step~3 is to prepare a clean, auditable, and minimally processed fNIRS dataset
for later analysis.
This step is intentionally simple and has three goals only:
\begin{enumerate}
    \item verify that each participant is linked to the correct fNIRS session folder and files;
    \item harmonize missing or incomplete triggers using the PsychoPy logs when justified;
    \item apply a basic preprocessing pipeline in MNE to produce files ready for Step~4.
\end{enumerate}

This step does \textbf{not} perform the final scientific analysis of the fNIRS signal.
Its function is to ensure that the data are complete enough, correctly matched, well documented,
and consistently preprocessed before any question-level or group-level interpretation.

\subsection{Core Design Choice}
The atomic unit of analysis in Step~3 is not a single file but a
\textbf{participant-session folder}, for example:
\begin{verbatim}
PID001/2025-11-06_008/
\end{verbatim}

All file checks, trigger checks, preprocessing, logging, and status assignment must be done
at the level of this participant-session folder.

\subsection{Raw Session Folder Structure}
A typical session folder contains files such as:
\begin{verbatim}
/PID001/2025-11-06_008/2025-11-06_008.nirs
/PID001/2025-11-06_008/2025-11-06_008.snirf
/PID001/2025-11-06_008/2025-11-06_008.snirf.bak_pre_marker_fix
/PID001/2025-11-06_008/2025-11-06_008.wl1
/PID001/2025-11-06_008/2025-11-06_008.wl2
/PID001/2025-11-06_008/2025-11-06_008.zip
/PID001/2025-11-06_008/2025-11-06_008_calibration.json
/PID001/2025-11-06_008/2025-11-06_008_config.hdr
/PID001/2025-11-06_008/2025-11-06_008_config.json
/PID001/2025-11-06_008/2025-11-06_008_description.json
/PID001/2025-11-06_008/2025-11-06_008_probeInfo.json
/PID001/2025-11-06_008/digpts.txt
\end{verbatim}

\subsection{Canonical File Usage in This Project}
To keep the pipeline simple and consistent, the files in each session folder should be used as follows:

\subsubsection*{Primary analysis file}
\begin{itemize}
    \item \texttt{<session>.snirf}
\end{itemize}

This is the main file to be loaded into MNE for trigger inspection, annotation handling,
quality control, and preprocessing.

\subsubsection*{Audit and traceability file}
\begin{itemize}
    \item \texttt{<session>.snirf.bak\_pre\_marker\_fix}
\end{itemize}

This file must be preserved as audit evidence.
If it exists, it should be compared once against the current \texttt{.snirf} to determine whether
the current file already contains previous marker edits.

\subsubsection*{Auxiliary metadata files}
\begin{itemize}
    \item \texttt{<session>\_probeInfo.json}
    \item \texttt{<session>\_config.json}
    \item \texttt{<session>\_description.json}
    \item \texttt{<session>\_calibration.json}
    \item \texttt{digpts.txt}
    \item \texttt{<session>\_config.hdr}
\end{itemize}

These files are not the main analysis input, but they should be used for metadata verification,
montage checks, and documentation.

\subsubsection*{Legacy or archival files}
\begin{itemize}
    \item \texttt{<session>.nirs}
    \item \texttt{<session>.wl1}
    \item \texttt{<session>.wl2}
    \item \texttt{<session>.zip}
\end{itemize}

These files should be retained for traceability, but the main Step~3 pipeline should not rely on them.

\subsection{Main Rule for Trigger Checking}
Trigger checking must be performed \textbf{within participant-session}.
Participants do not all need to have the same total number of triggers, because some participants
did not finish the full task.
A participant-session should be considered valid when:
\begin{itemize}
    \item the PsychoPy log indicates an early stop or incomplete session, and
    \item the fNIRS trigger sequence matches that participant's actual completed portion of the task,
    either directly or after justified trigger repair.
\end{itemize}

Therefore, Step~3 compares:
\begin{quote}
fNIRS triggers for one participant-session \quad versus \quad PsychoPy-derived expected triggers
for the same participant-session.
\end{quote}

\subsection{Scope}
Step~3 includes:
\begin{itemize}
    \item session inventory and validation of raw files;
    \item extraction of native triggers from the \texttt{.snirf} file;
    \item reconstruction of expected triggers from PsychoPy;
    \item trigger comparison and repair when missing triggers can be safely inserted;
    \item participant-session quality control;
    \item basic preprocessing of the fNIRS signal;
    \item full documentation of all issues and final session status.
\end{itemize}

Step~3 excludes:
\begin{itemize}
    \item condition-level fNIRS analysis;
    \item statistical comparison between abstract and concrete questions;
    \item ROI-level scientific interpretation;
    \item final figures and inferential results.
\end{itemize}

\subsection{Guiding Principles}
\begin{enumerate}
    \item Keep the workflow simple.
    \item Never overwrite the raw session files.
    \item Use the \texttt{.snirf} file as the working raw file.
    \item Preserve the provenance of every trigger and every repair.
    \item Do not silently exclude problematic sessions.
    \item Log all participant-session issues explicitly.
    \item Use a single standard preprocessing path unless a session clearly fails.
\end{enumerate}

\subsection{Inputs}
Step~3 requires:
\begin{itemize}
    \item the participant-session folder for each PID;
    \item the corresponding PsychoPy task log;
    \item the clean Step~2 question reconstruction, or the PsychoPy event reconstruction used in Step~2;
    \item the final montage information for the study;
    \item the information that the final montage is left-lateralized and includes short channels and accelerometer data.
\end{itemize}

\subsection{Expected Outputs}

\subsubsection*{Manifest and trigger tables}
\begin{verbatim}
01_fnirs_session_manifest.csv
02_fnirs_trigger_raw.csv
03_psychopy_trigger_expected.csv
04_trigger_comparison_by_session.csv
05_trigger_repair_log.csv
06_fnirs_trigger_harmonized.csv
\end{verbatim}

\subsubsection*{Quality-control tables}
\begin{verbatim}
07_fnirs_session_status.csv
08_fnirs_channel_qc.csv
09_fnirs_problem_log.csv
10_fnirs_file_audit.csv
\end{verbatim}

\subsubsection*{Preprocessing outputs}
\begin{verbatim}
11_preprocessing_parameters.json
12_preprocessed_file_manifest.csv
data_clean/step3/preprocessed/PID001_2025-11-06_008_preprocessed_raw.fif
data_clean/step3/preprocessed/PID002_2025-11-06_009_preprocessed_raw.fif
...
\end{verbatim}

\subsubsection*{Reports}
\begin{verbatim}
reports/step3/step3_fnirs_preprocessing_report.tex
reports/step3/session_notes/PID001_2025-11-06_008.tex
reports/step3/session_notes/PID002_2025-11-06_009.tex
...
\end{verbatim}

\subsection{Internal Structure of Step~3}
To keep the workflow very clear, Step~3 is divided into four blocks:
\begin{enumerate}
    \item \textbf{Step~3A: Session inventory and file validation}
    \item \textbf{Step~3B: Trigger extraction and harmonization}
    \item \textbf{Step~3C: Session-level quality control}
    \item \textbf{Step~3D: Basic preprocessing in MNE}
\end{enumerate}

\subsection{Step~3A: Session Inventory and File Validation}

\subsubsection{Step~3A.1: Build the session manifest}
Create one row per participant-session folder.
For each row, record:
\begin{itemize}
    \item \texttt{participant\_id}
    \item \texttt{session\_id}
    \item session folder path
    \item \texttt{snirf\_file}
    \item \texttt{snirf\_backup\_file}
    \item \texttt{psychopy\_file}
    \item file timestamp
    \item whether the session folder is readable
    \item notes
\end{itemize}

\subsubsection{Step~3A.2: Verify expected file presence}
For each session folder, verify whether the following exist:
\begin{itemize}
    \item main \texttt{.snirf} file;
    \item \texttt{.snirf.bak\_pre\_marker\_fix} file, if available;
    \item \texttt{\_probeInfo.json};
    \item \texttt{\_config.json};
    \item \texttt{\_description.json};
    \item \texttt{\_calibration.json};
    \item \texttt{digpts.txt};
    \item paired PsychoPy file.
\end{itemize}

The absence of some metadata files may not exclude the session, but it must be documented.

\subsubsection{Step~3A.3: Define analysis and audit roles for files}
Each session manifest row should explicitly record:
\begin{itemize}
    \item which file is the primary raw analysis file;
    \item which file is the audit backup;
    \item which files are auxiliary metadata only;
    \item whether any previous marker-fix backup exists.
\end{itemize}

\subsubsection{Step~3A.4: Check raw file identity}
Before any preprocessing, confirm that the \texttt{.snirf} file loaded for one session truly belongs
to that participant-session by checking:
\begin{itemize}
    \item session folder name;
    \item file stem;
    \item session timestamp;
    \item correspondence to the PsychoPy session.
\end{itemize}

\subsection{Step~3B: Trigger Extraction and Harmonization}

\subsubsection{Step~3B.1: Extract native triggers from the \texttt{.snirf} file}
For each session, extract the native events or annotations present in the current \texttt{.snirf}.
Store at least:
\begin{itemize}
    \item \texttt{participant\_id}
    \item \texttt{session\_id}
    \item \texttt{event\_time\_fnirs}
    \item \texttt{event\_label\_raw}
    \item \texttt{event\_code\_raw}
    \item \texttt{source = snirf}
\end{itemize}

\subsubsection{Step~3B.2: Audit the backup marker file when present}
If \texttt{.snirf.bak\_pre\_marker\_fix} exists, compare its event table with the current \texttt{.snirf}.
The goal is not to create a second pipeline, but simply to document:
\begin{itemize}
    \item whether the current \texttt{.snirf} already differs from the backup;
    \item whether previous manual or scripted marker fixes likely occurred;
    \item whether the current file should be treated as already modified.
\end{itemize}

\subsubsection{Step~3B.3: Build expected triggers from PsychoPy}
From the PsychoPy log, reconstruct the expected sequence of task events for the same participant-session.
At minimum, the expected trigger table should include:
\begin{itemize}
    \item \texttt{participant\_id}
    \item \texttt{session\_id}
    \item \texttt{event\_order}
    \item \texttt{event\_type}
    \item \texttt{event\_time\_psychopy}
    \item \texttt{block}
    \item \texttt{trial\_idx\_in\_block}
    \item \texttt{question\_id}
    \item \texttt{source = psychopy}
\end{itemize}

\subsubsection{Step~3B.4: Define the minimum harmonized trigger set}
To keep Step~3 simple, the minimum trigger set for harmonization should be:
\begin{itemize}
    \item \texttt{block\_start}
    \item \texttt{question\_start}
    \item \texttt{answer}
    \item \texttt{block\_end}
\end{itemize}

For this project, \texttt{question\_start} should be defined as the PsychoPy event corresponding to the
beginning of question presentation.

\subsubsection{Step~3B.5: Compare fNIRS and PsychoPy trigger sequences}
For each participant-session, compare:
\begin{itemize}
    \item total number of question-start events;
    \item total number of answer events;
    \item total number of block-start events;
    \item total number of block-end events;
    \item event ordering.
\end{itemize}

A session should then be assigned one comparison result:
\begin{itemize}
    \item \texttt{MATCH\_RAW}
    \item \texttt{MATCH\_INCOMPLETE\_SESSION}
    \item \texttt{MISSING\_TRIGGERS\_REPAIRABLE}
    \item \texttt{MISMATCH\_REVIEW\_REQUIRED}
\end{itemize}

\subsubsection{Step~3B.6: Repair missing triggers when justified}
If the fNIRS file is missing triggers but the alignment with PsychoPy is still reliable,
create a harmonized trigger table.
The repair logic must remain simple:
\begin{enumerate}
    \item keep all original fNIRS triggers;
    \item add only the missing triggers from PsychoPy;
    \item never delete original trigger information from the audit tables;
    \item mark every inserted trigger explicitly as repaired;
    \item do not force repair when alignment is uncertain.
\end{enumerate}

\subsubsection{Step~3B.7: Store trigger provenance}
Every harmonized trigger row should contain:
\begin{itemize}
    \item \texttt{participant\_id}
    \item \texttt{session\_id}
    \item \texttt{event\_type}
    \item \texttt{event\_time\_final}
    \item \texttt{event\_source\_final}
    \item \texttt{raw\_snirf\_present}
    \item \texttt{inserted\_from\_psychopy}
    \item \texttt{repair\_confidence}
    \item \texttt{notes}
\end{itemize}

\subsection{Step~3C: Session-Level Quality Control}

\subsubsection{Step~3C.1: Verify recording readability and duration}
For each participant-session, confirm:
\begin{itemize}
    \item the \texttt{.snirf} file opens successfully in MNE;
    \item the sampling information is readable;
    \item the recording duration is non-zero and plausible;
    \item the annotation structure is readable.
\end{itemize}

\subsubsection{Step~3C.2: Verify channel presence}
For each participant-session, record:
\begin{itemize}
    \item total number of channels;
    \item number of long channels;
    \item number of short channels;
    \item presence or absence of accelerometer information;
    \item whether channel names are complete and consistent.
\end{itemize}

\subsubsection{Step~3C.3: Verify montage-related metadata}
Using \texttt{\_probeInfo.json}, \texttt{digpts.txt}, and related metadata files, verify:
\begin{itemize}
    \item that the expected montage information exists;
    \item that the session appears to correspond to the intended left-sided montage;
    \item that short-channel information is available when expected;
    \item that there are no obvious file-level inconsistencies.
\end{itemize}

\subsubsection{Step~3C.4: Log participant-session problems}
Every session with any issue must be written to a problem log.
Typical issues include:
\begin{itemize}
    \item missing \texttt{.snirf} file;
    \item unreadable \texttt{.snirf} file;
    \item unresolved trigger mismatch;
    \item missing or unclear short channels;
    \item missing accelerometer information;
    \item implausibly short recording;
    \item preprocessing failure.
\end{itemize}

\subsubsection{Step~3C.5: Assign a session QC label}
Each participant-session should receive one quality-control label:
\begin{itemize}
    \item \texttt{OK}
    \item \texttt{OK\_WITH\_WARNINGS}
    \item \texttt{REVIEW\_REQUIRED}
    \item \texttt{EXCLUDE}
\end{itemize}

\subsection{Step~3D: Basic Preprocessing in MNE}

\subsubsection{Step~3D.1: Load the main \texttt{.snirf} file}
Use the main \texttt{.snirf} file as the raw analysis file for each participant-session.
At this stage, attach or verify:
\begin{itemize}
    \item participant ID;
    \item session ID;
    \item harmonized annotations or equivalent event structure;
    \item bad-channel markings when identified.
\end{itemize}

\subsubsection{Step~3D.2: Convert raw intensity to optical density}
This should be the first signal-processing step in the pipeline.

\subsubsection{Step~3D.3: Compute basic channel quality}
On optical density data, compute a simple channel quality metric and document poor channels.
Channels marked as problematic must remain visible in the QC tables and reports.

\subsubsection{Step~3D.4: Apply one motion-correction step}
To keep the pipeline basic, use a single motion-correction method.
A simple default is to apply TDDR before conversion to haemoglobin.

\subsubsection{Step~3D.5: Convert to haemoglobin concentration}
After the optical-density stage and motion correction, convert the data to haemoglobin concentration.

\subsubsection{Step~3D.6: Apply simple filtering}
Apply one standard filter configuration to the haemoglobin data.
The exact parameters used must be stored in:
\begin{verbatim}
11_preprocessing_parameters.json
\end{verbatim}

\subsubsection{Step~3D.7: Use short channels when available}
If short channels are present and valid, apply short-channel correction in the standard branch.
If they are missing or not usable, document this clearly and continue with the same basic pipeline
without adding unnecessary complexity.

\subsubsection{Step~3D.8: Retain accelerometer information for QC}
If accelerometer data are present in the session, retain them for inspection and documentation.
At this stage, accelerometer data do not need a separate advanced modeling branch.

\subsubsection{Step~3D.9: Save the preprocessed output}
For each participant-session, save one preprocessed output file that will be used in Step~4.
The saved file must contain:
\begin{itemize}
    \item preprocessed fNIRS signal;
    \item harmonized trigger annotations or equivalent event structure;
    \item marked bad channels;
    \item participant-session identifiers.
\end{itemize}

\subsection{Documentation Required for Every Participant-Session}
For complete traceability, the report should state for every participant-session:
\begin{itemize}
    \item whether the main \texttt{.snirf} file was found;
    \item whether a backup pre-marker-fix \texttt{.snirf} file exists;
    \item whether the current \texttt{.snirf} appears already marker-modified;
    \item whether the paired PsychoPy file was found;
    \item whether the trigger sequence matched directly;
    \item whether trigger repair was needed;
    \item whether the recording was complete or incomplete;
    \item whether short channels were present;
    \item whether accelerometer information was present;
    \item how many channels were marked problematic;
    \item whether preprocessing completed successfully;
    \item final session status.
\end{itemize}

\subsection{Quality-Control Checks}
Before Step~3 is considered complete, the following checks should be run:
\begin{enumerate}
    \item every participant-session folder has a manifest entry;
    \item every session has one declared primary raw file;
    \item every session has one trigger-comparison result;
    \item repaired triggers are explicitly marked;
    \item unresolved trigger mismatches are logged;
    \item the presence of short channels is recorded;
    \item the presence of accelerometer information is recorded;
    \item bad channels are documented;
    \item preprocessing parameters are saved;
    \item every output file is linked to exactly one participant-session;
    \item every participant-session has a final status.
\end{enumerate}

\subsection{Final Status Labels}
At the end of Step~3, each participant-session should receive one final status:
\begin{itemize}
    \item \texttt{READY\_FOR\_STEP4}
    \item \texttt{READY\_FOR\_STEP4\_WITH\_WARNINGS}
    \item \texttt{REVIEW\_BEFORE\_STEP4}
    \item \texttt{EXCLUDE\_FROM\_STEP4}
\end{itemize}

\subsection{Minimal Completion Criteria for Step~3}
Step~3 is complete when:
\begin{itemize}
    \item all participant-session folders have been inventoried;
    \item all \texttt{.snirf} files have been checked;
    \item all expected PsychoPy triggers have been reconstructed;
    \item trigger comparison has been completed for every session;
    \item missing triggers have been repaired when justified;
    \item unresolved sessions are explicitly documented;
    \item a harmonized trigger table exists;
    \item a basic preprocessing pipeline has been completed;
    \item preprocessed output files exist for all usable sessions;
    \item participant-session QC and problem logs are complete.
\end{itemize}

\subsection{Handoff to Step~4}
Step~4 should start only from the outputs of Step~3:
\begin{itemize}
    \item preprocessed fNIRS files;
    \item harmonized trigger table;
    \item session-level QC table;
    \item problem log;
    \item preprocessing parameter file.
\end{itemize}

\subsection{Short Practical Summary}
In simple terms, Step~3 should:
\begin{quote}
use the participant-session folder as the working unit, load the \texttt{.snirf} file,
check and repair triggers using PsychoPy when needed, verify session quality, run a basic
MNE preprocessing pipeline, and document every issue before passing the data to Step~4.
\end{quote}

\end{document}
